\section{Amenability from random polygonal partitions}\label{sec:amenability}

The input to this section is the polygonal Boolean algebra and its translation
arithmetic from Lemma~\ref{lem:arithmetic}.  We prove amenability for every
fixed finite number of tracks.  The main construction is a random finite
polygonal partition: with high probability its cells lie inside the translation
charts of prescribed exchanges, and its law is almost unchanged when they move the
cells.  We state the precise coupling target below, before constructing the
random fields that will supply its partition edges.

The use of almost-invariant finite configurations is related to
\cite[Theorem~C]{JuschenkoMonod2013} and the extensive-amenability
methods of \cite{JuschenkoMatteBonMonodDeLaSalle2018}.  The proof here
constructs the correlated partition law and its conversion to group
measures directly.  This is essential in dimension two, where even
free minimal Cantor $\ZZ^2$-actions can have nonamenable full
groups \cite[Theorem~1]{ElekMonod2013}.

\begin{theorem}\label{thm:amenability}
For every positive integer $m$, the discrete group $F_m$ is amenable.  Every
subgroup of $F_m$, in particular $A_m$, satisfies the finite-set F\o lner
condition.
\end{theorem}

Throughout the proof $m$ and a finite set $\mathcal K\subset F_m$ are fixed.
Include inverses in $\mathcal K$ if necessary.  Choose finite polygonal
continuity partitions for its elements, so that each piece is translated by
one vector in $\OO^2$ to a specified track.  Subdivide further along supporting
lines to obtain convex pieces.  All source and image boundaries are contained
in a fixed finite collection of allowable line segments.  Integer translations
used to cross the edges of the square are included among the translation
vectors.  Consequently their coordinates and conjugate coordinates, and the
levels and conjugate levels of all these supporting lines, are bounded by a
constant depending only on $\mathcal K$.

For each $\delta>0$ we will construct a law on finite polygonal partitions
with the following property.  For every $g\in\mathcal K$, two partitions
$\mathcal P,\mathcal P'$ with this law can be coupled so that, with
probability greater than $1-\delta$, every cell of $\mathcal P$ lies in
one translation chart of $g$ and
\[
 \mathcal P'=g\mathcal P=\{gC:C\in\mathcal P\}.
\]
Each cell is contained in a single track.  Two cells have the same shape
if they differ by an ordinary $\OO^2$ translation, ignoring their tracks.
On the displayed event, $g$ preserves the multiset of cell shapes.
For each occurring multiset, fix a baseline partition.  Uniformly matching
its cells to the sampled cells of the same shape will lift this coupling
to almost-invariant probability measures on $F_m$.  The last subsection
proves that lift and its conversion to finite F\o lner sets.

\subsection{Lattice counts and sites with a specified side}

Let $D$ be as in Lemma~\ref{lem:arithmetic}, and set
\[
 L=D^{-1}\OO^2,
 \qquad
 \Lambda=\{(z,\bar z):z\in L\}\subset\RR^2\times\RR^2.
\]
Write $\rho$ for the density of this four-dimensional lattice.  Norms on the
ordinary and conjugate planes may be taken to be the maximum norm; replacing
them by Euclidean norms changes only constants.

We record the quantitative consequences of unit rescaling that will be used.
For intervals $I,J\subset\RR$ of lengths $a,b>0$,
\begin{equation}\label{eq:am-one-count}
 \#\{u\in D^{-1}\OO:u\in I,\ \bar u\in J\}\le C(ab+1).
\end{equation}
To prove this when $ab\ge1$, rescale by a power of the unit $\tau$ so that
the two side lengths become comparable to $\sqrt{ab}$.  A fixed fundamental
parallelogram then gives the bound by counting tiles meeting the rectangle.
When $ab<1$, either the same argument applies after enlarging the rectangle
to area one, or one uses
$|u\bar u|\ge D^{-2}$ for a nonzero difference $u$ and subdivides the
rectangle into a fixed number of smaller rectangles.  This also proves that
the constant is independent of the positions of $I,J$.

There is a corresponding uniform Riemann-sum statement.  For ordinary scale
$s>0$ and conjugate scale $t>0$, the lattice in coordinates $(z/s,\bar z/t)$
has fundamental tiles of diameter at most
\begin{equation}\label{eq:am-tile}
 C(st)^{-1/2}.
\end{equation}
Indeed, apply a unit with ordinary size comparable to $\sqrt{s/t}$ to a
fixed basis of $\Lambda$ before scaling the two coordinates.  Both parts of
each resulting basis vector have size $O((st)^{-1/2})$.  The tile volume is
$1/(\rho s^2t^2)$.  If $st\to\infty$, integrating a compactly supported
Lipschitz function over these tiles and replacing its value by the value at
the lattice point changes the normalized sum by
$O((st)^{-1/2})$, uniformly over translates of the lattice.  Tiles meeting
the boundary of a fixed box contribute the same order of error.  Thus, in
particular, a box of ordinary side lengths $O(s)$ and conjugate side lengths
$O(t)$ contains $O(s^2t^2)$ sites once $st\ge1$.

At a point on a cut, a side must be specified.  Fix vectors $v,w\in\RR^2$
such that $v\ne0$ and $\ell_j(w)\ne0$ whenever $\ell_j(v)=0$.  The
\emph{flag} $(v,w)$ assigns the sign of $\ell_j-c$ at $z$ to be the
lexicographic sign of
\[
 (\ell_j(z)-c,\ell_j(v),\ell_j(w)).
\]
This is a formal convention, denoted $z+\varepsilon v+\varepsilon^2w$;
no common numerical choice of $\varepsilon$ for the infinitely many cuts is
intended.  Each Boolean polygon has a well-defined value at this flag.
On the torus, represent a flagged point by the unique lift in the closed
square whose flag enters the square.  For a fixed flag type, let $\mathcal S$
be the countable set of these representatives in the $m$ tracks with base
point in $L$ modulo $\ZZ^2$.  Each exchange permutes $\mathcal S$: use the
translation belonging to the Boolean chart selected by the flag, and then
choose the square representative again.  The inverse exchange proves
bijectivity.  Translation leaves the two flag vectors unchanged.

All subsequent field estimates hold for any one fixed flag type, with
constants allowed to depend on that type.  A finite number of independent
copies, also with different flag types, will be used for the four directions.
Boundary representatives affect none of the counts: along a fixed allowable
line in a bounded ordinary region and a conjugate ball of radius $N$, there
are $O(N)$ sites, by~\eqref{eq:am-one-count} in a tangent coordinate.

The partitions will use two bit fields for each of the four line directions.
On a finite collection of allowable line segments in the squares, the first
field marks subdivision points.  The second field decides whether the
interval beginning at each marked point is a barrier.  Together with the
square edges, these barriers cut each track into polygonal cells.  An
interval uses the second bit at its initial endpoint, interpreted by a flag
pointing along the interval into its chart.  Its terminal endpoint will be
forced independently when it lies on a chart edge.  This is why we use
tangent flags.  The exact rule, including the square endpoints, and its
covariance verification follow the field estimates.

The joint law of these fields must change little under each exchange in
$\mathcal K$.  At the same time, they must force prescribed bits at all
relevant sites: chart crossings are subdivision points, chart edges are
complete barriers, and sufficiently distant sites in the conjugate plane
are unmarked.  We next construct general fields with these two properties;
the specific fields for the partition will then be chosen from the fixed
chart data.

\subsection{Positive averaging matrices}

Let $\theta:\RR^2\to\RR$ be smooth and constant outside a compact set.
Choose $\chi\in C_c^\infty(\RR^2)$ with $0\le\chi\le1$ that equals one
on a fixed neighborhood of $\supp\nabla\theta$.  We can enlarge this
neighborhood whenever finitely many signals are being treated.  The mean
field is
\[
 \mu_N(z)=\theta(\bar z/N),
\]
with the same signal on every track.  Fix small positive parameters
$\eta,\kappa$, with $4\eta$ less than the width of the neighborhood where
$\chi=1$.  Eventually $\eta,\kappa$ will be chosen first and $N$ will then
tend to infinity.

We will construct a positive semidefinite matrix $B_N$ with finitely many
nonzero rows and columns.  We add bounded independent centered noise $\xi_z$
and a correlated perturbation $B_N\xi$ to the mean, then take the sign of
the resulting field.  A bounded chart translation changes $\mu_N$ by $O(N^{-1})$ on
$O(N^2)$ sites, so the $\ell^2$ norm of the mean difference need not tend
to zero.  The matrix $B_N$ is designed to multiply
this slowly varying difference approximately by a factor tending to
infinity.  Applying $(I+B_N)^{-1}$ will therefore make the mean difference
small in $\ell^2$, as needed for the total-variation estimate below.
Meanwhile, the $\ell^2$ norms of the rows of $B_N$ will be
sufficiently small that $\|B_N\xi\|_\infty\le1/2$ with probability tending
to one.  These are compatible demands: averaging preserves a slowly
varying signal while cancelling independent centered noise.

Let $\mathcal R_N$ be the powers of two between $N^{1/4}$ and $N^{1/2}$.
For $r\in\mathcal R_N$ put $s=r/N$ and
\[
 K(u)=(1-|u_1|)_+(1-|u_2|)_+.
\]
The lines in the next definition are auxiliary separators for the averaging
matrix.  We average over their thresholds to obtain a deterministic matrix;
the partition barriers will instead be sampled from the bit fields.  The
auxiliary separators suppress matrix entries across chart edges while
retaining averaging at most sites.

Independently for each of the four direction families choose a threshold
uniformly in $[\kappa/s,2\kappa/s]$, and cut the plane along all lines
$\ell_j=c$, $c\in\OO$, with $|\bar c|$ at most that threshold.
There are only finitely many such lines in any bounded ordinary region,
by~\eqref{eq:am-one-count}.  The flags determine on which side a site lies.
Let $P_r(z,w)$ be the probability that no chosen line separates the two
flagged sites.  This definition is used for sites on the same track.

Define a matrix, zero between different tracks, by
\begin{equation}\label{eq:am-matrix}
 M_r(z,w)=
 \frac{\chi(\bar z/N)\chi(\bar w/N)}{\rho(r\eta)^2}
 K\left(\frac{z-w}{s}\right)
 K\left(\frac{\bar z-\bar w}{N\eta}\right)P_r(z,w).
\end{equation}
Only $O(N^2)$ rows or columns can be nonzero.  This follows by applying the
box count to the bounded ordinary square and a conjugate ball of radius
$O(N)$.  Every $M_r$ is positive semidefinite.  Indeed,
$(1-|u|)_+$ is the overlap length of two translates of a unit interval, so
its convolution kernel is positive semidefinite.  Products give the two
kernels in~\eqref{eq:am-matrix}.  For a fixed collection of separating lines,
the matrix that is one exactly when two sites lie in the same chamber is
positive semidefinite, since its quadratic form is the sum, over chambers,
of the squares of the coordinate sums.  Averaging gives positivity of
$P_r$.  Entrywise products preserve positivity: Gram representations turn
them into inner products of tensor products.  Finally multiplication by
$\chi$ on the two sides and the direct sum over tracks preserve positivity.

We use the finite matrix
\begin{equation}\label{eq:am-B}
 B_N=(\log N)^{-3/4}\sum_{r\in\mathcal R_N}M_r,
 \qquad
 W_N=(\log N)^{-3/4}|\mathcal R_N|.
\end{equation}
In particular $W_N\asymp(\log N)^{1/4}\to\infty$ and $I+B_N$ is
positive definite with inverse of operator norm at most one.

\begin{lemma}\label{lem:am-estimates}
For each $g\in\mathcal K$, let $U_g$ be its permutation operator on
$\ell^2(\mathcal S)$ and put
$h_N=U_g\mu_N-\mu_N$.  With $\eta,\kappa$ fixed,
\begin{align}
 \|U_gB_NU_g^{-1}-B_N\|_{\HS}&\longrightarrow0,
       \label{eq:am-covariance}\\
 \Pr\{\|B_N\xi\|_\infty>1/2\}&\longrightarrow0,
       \label{eq:am-concentration}
\end{align}
where the $\xi_z$ are independent, centered random variables supported on
$[-1,1]$.  Moreover
\begin{equation}\label{eq:am-signal-bound}
 \limsup_{N\to\infty}\|(I+B_N)^{-1}h_N\|_2
       \le C(\eta+\sqrt\kappa),
\end{equation}
where $C$ is independent of sufficiently small $\eta,\kappa$.  All assertions
hold uniformly over the fixed finite set $\mathcal K$.
\end{lemma}

\begin{proof}
We give separately the covariance, concentration, and signal estimates.
The count following~\eqref{eq:am-tile}, now with $t=N\eta$, shows that a
row of $M_r$ has $O_\eta(r^2)$ nonzero entries, each
$O_\eta(r^{-2})$.  Therefore
\begin{equation}\label{eq:am-row-norms}
 \|M_r(z,\cdot)\|_2\le C_\eta/r.
\end{equation}
The row sums are bounded by an absolute constant for all sufficiently large
$N$ at any fixed $\eta$: omit $P_r$ and $\chi$, and use the uniform
Riemann sum for $K(u)K(v)$, whose integral is one.  In particular this bound
can be chosen independently of small $\eta,\kappa$ once $N$ is large enough.

Write $\Delta_r=U_gM_rU_g^{-1}-M_r$.  Compare its entries by comparing
$(z,w)$ with $(gz,gw)$.  If $z,w$ lie in different continuity pieces on the
same source track, some supporting line of the fixed continuity partition
separates their flags.  To see this, assign to a point the signs of all the
finitely many supporting lines; each sign chamber lies in one continuity
piece.  Thus different pieces have different sign assignments.  Its
conjugate level is bounded independently of $N$, so it is included for every
threshold $\kappa/s$ when $N$ is large.  The original separator factor is
zero.  Their images lie in different image pieces; the same reasoning makes
the image factor zero if they are on the same track, and otherwise the
image matrix entry is zero by definition.  This also deals with points
on different source tracks whose images share a track.  Hence only pairs
in a common continuity piece need be considered.

For such a pair, $gz=z+u$ and $gw=w+u$ in the chosen square charts, for a
fixed $u\in\OO^2$.  Both difference kernels in~\eqref{eq:am-matrix} are
unchanged.  In one direction let $d$ be the least absolute conjugate level
of a line separating $z,w$, taking $d=+\infty$ if there is none.  The
probability that the threshold is less than $d$ is a function of $d$ with
Lipschitz constant $s/\kappa$.  Translation bijects the separating lines,
changing their conjugate levels by $\overline{\ell_j(u)}$.  Either both minima are infinite,
or they differ by a bounded amount.  Taking the product over the four
directions gives a change $O_\kappa(s)$ in $P_r$.  Each cutoff factor changes
by $O(N^{-1})$.  We have proved, on the union of the two entry supports,
\begin{equation}\label{eq:am-entry-change}
 |\Delta_r(z,w)|\le \frac{C_{\eta,\kappa}}{Nr}.
\end{equation}
This union contains $O_\eta(N^2r^2)$ pairs, also after permutation.
For $r\le r'$, the intersection of the two entry supports contains at most
the number of pairs in the smaller support.  It follows that
\[
 |\langle\Delta_r,\Delta_{r'}\rangle_{\HS}|
 \le C_{\eta,\kappa}\frac{r}{r'}.
\]
The scales are dyadic, so the sum of $r/r'$ over $r\le r'$ is
$O(|\mathcal R_N|)$.  Consequently
\[
 \|U_gB_NU_g^{-1}-B_N\|_{\HS}^2
 \le C_{\eta,\kappa}(\log N)^{-3/2}|\mathcal R_N|
 =O_{\eta,\kappa}((\log N)^{-1/2}),
\]
which proves~\eqref{eq:am-covariance}.

For concentration,~\eqref{eq:am-row-norms} and the geometric sum over scales
give
\[
 \|B_N(z,\cdot)\|_2
 \le C_\eta(\log N)^{-3/4}N^{-1/4}.
\]
The rowwise bounded-independent-summand estimate is Hoeffding's
inequality~\cite[Theorem~2]{Hoeffding1963}, rederived here.
If $\xi$ is centered and $|\xi|\le1$, convexity bounds its exponential
moment by $\cosh t\le e^{t^2/2}$.  Multiplication of the independent
exponential moments and optimization in $t$ therefore give
\[
 \Pr\{|(B_N\xi)_z|>1/2\}
 \le 2\exp\{-c_\eta N^{1/2}(\log N)^{3/2}\}.
\]
There are $O(N^2)$ active rows; outside them $B_N\xi=0$.  The union bound
proves~\eqref{eq:am-concentration}.

It remains to explain why the matrices average the signal difference.
On an image chart whose source translation is $u$, the exact formula is
\[
 h_N(z)=\theta\left(\frac{\bar z-\bar u}{N}\right)
              -\theta\left(\frac{\bar z}{N}\right).
\]
Thus $|h_N|\le C/N$, it is supported on $O(N^2)$ sites, and $Nh_N$,
viewed as a function of $\bar z/N$, has a uniformly bounded Lipschitz
constant on each fixed chart.  In particular $\|h_N\|_2\le C$.
For large $N$ its support, together with its $2\eta$-neighborhood in scaled
conjugate coordinates, lies where $\chi=1$ whenever $\eta$ is sufficiently
small.  This guarantees that the cutoff introduces no error in averaging
$h_N$.

We bound the exceptional rows where a separator or a chart boundary can
interfere with this averaging.  For each direction $\ell_j$ one may complete
it to an $\OO$-unimodular coordinate system, using respectively
\[
 (x,y),\quad(y,x),\quad(y-\lambda x,x),\quad(x-\lambda y,y).
\]
In these coordinates the paired lattice is again a product of two
one-dimensional quadratic lattices.  In a bounded ordinary region and a
conjugate ball of radius $O(N)$, the number of sites within ordinary distance
$C s$ of a line $\ell_j=c$ is, by~\eqref{eq:am-one-count}, at most
\[
 C(sN+1)(N+1)\le C sN^2,
\]
since $sN=r\to\infty$.  The number of lines meeting that ordinary region
with $|\bar c|\le2\kappa/s$ is $O(\kappa/s+1)$.  The rows within distance
$Cs$ of any of them consequently number at most
\begin{equation}\label{eq:am-bad-rows}
 C(\kappa+s)N^2.
\end{equation}
The same estimate with only finitely many lines gives $O(sN^2)$ rows near
the fixed chart boundaries and square edges.  These estimates include rows
exactly on a line, independently of the flag convention.

Outside these exceptional rows, every contributing neighbor is in the same
ordinary chart, and no line allowed in the definition of $P_r$ separates it
from the row.  Hence $P_r=1$.  The normalized Riemann-sum argument
in~\eqref{eq:am-tile}, together with the Lipschitz bound on $Nh_N$, gives
\[
 |(M_rh_N)(z)-h_N(z)|
 \le \frac{C}{N}\left(\eta+(r\eta)^{-1/2}\right).
\]
The assertion also holds for rows where $h_N(z)=0$: a contributing nonzero
value of $h_N$ still lies in the cutoff-one neighborhood.  On exceptional
rows the bound $C/N$ suffices, since the row sums are bounded.  Outside a
conjugate ball of radius $RN$, with $R$ fixed, both terms vanish.  Squaring and summing yields,
uniformly over $r\in\mathcal R_N$,
\begin{equation}\label{eq:am-approx-identity}
 \|M_rh_N-h_N\|_2
 \le C\left(\eta+\sqrt\kappa+\sqrt{s}+(r\eta)^{-1/2}\right)
 = C(\eta+\sqrt\kappa)+o(1).
\end{equation}
The constant $C$ in this estimate is independent of sufficiently small
$\eta,\kappa$; the $o(1)$ is taken with them fixed.

Finally put $e_N=B_Nh_N-W_Nh_N$.  Equation~\eqref{eq:am-approx-identity}
implies
$\|e_N\|_2\le W_N(C(\eta+\sqrt\kappa)+o(1))$.  The identity
\[
 (I+B_N)h_N=(1+W_N)h_N+e_N
\]
and the contraction bound for $(I+B_N)^{-1}$ give
\[
 \|(I+B_N)^{-1}h_N\|_2
 \le \frac{\|h_N\|_2}{1+W_N}
   +\frac{W_N}{1+W_N}\bigl(C(\eta+\sqrt\kappa)+o(1)\bigr).
\]
Since $W_N\to\infty$, this proves~\eqref{eq:am-signal-bound}.
\end{proof}

The three estimates have different purposes.  The matrix covariance controls
a change of the noise law.  The inverse signal estimate makes a bounded
Euclidean displacement of the mean inexpensive in that law.  Concentration
will force marks on specified regions, simultaneously at every relevant
site.  We now put these statements together in total variation on the whole
countable set of sites.

\subsection{Total variation of the entire field}

Choose a nonnegative even smooth function $f$ supported on $[-1,1]$, positive
on $(-1,1)$, with $\int f^2=1$.  Give each $\xi_z$ the density $f^2$ and
make the coordinates independent.  Define
\begin{equation}\label{eq:am-field}
 Y_N=\mu_N+(I+B_N)\xi,
 \qquad b_N(z)=\mathbf1_{\{Y_N(z)>0\}}.
\end{equation}
The next elementary finite-dimensional calculation supplies a bound whose
constant does not grow with the number of sites.

\begin{lemma}\label{lem:am-density-speed}
Let $\xi\in\RR^d$ have independent coordinates of density $f^2$.  Suppose
$T_t$ is a differentiable path of invertible real matrices of positive
determinant and $\mu_t\in\RR^d$ is differentiable.  If $p_t$ is the density
of $\mu_t+T_t\xi$, then
\[
 \left\|\frac{d}{dt}\sqrt{p_t}\right\|_{L^2}
 \le C_f\left(\|T_t^{-1}\dot T_t\|_{\HS}
                  +\|T_t^{-1}\dot\mu_t\|_2\right),
\]
with $C_f$ independent of $d$.
\end{lemma}

\begin{proof}
Set $F(x)=\prod_{i=1}^d f(x_i)$,
$L=T_t^{-1}\dot T_t$ and $v=T_t^{-1}\dot\mu_t$.
Differentiate
$(\det T_t)^{-1/2}F(T_t^{-1}(y-\mu_t))$ and pull back by $y=\mu_t+T_tx$.
The $L^2$ norm of the derivative is the norm of
\[
 (Lx+v)\cdot\nabla F+\tfrac12\operatorname{tr}(L)F.
\]
Under the product probability density $F^2$, write
$a(x)=f'(x)/f(x)$ on the interior of the support.  This notation is
legitimate in $L^2(f^2\,dx)$ because $\int(f')^2<\infty$.
Integration by parts gives
\[
 \mathbb E a(\xi_i)=0,\quad
 \mathbb E[\xi_i a(\xi_i)]=-\tfrac12,\quad
 \mathbb E\xi_i=0.
\]
The diagonal summands are
$L_{ii}(\xi_i a(\xi_i)+1/2)$ and the mean summands are $v_i a(\xi_i)$.
They have mean zero, uniformly bounded second moments, and different
coordinates are orthogonal.  The two types on the same coordinate are
orthogonal as well, by parity.  An off-diagonal summand is
$L_{ij}\xi_j a(\xi_i)$; it has mean zero separately in each of its two
coordinates.  Therefore it is orthogonal to every single-coordinate
summand and to off-diagonal summands with a different unordered coordinate
pair.  The two terms for $(i,j)$ and $(j,i)$ need not be orthogonal, but
Cauchy--Schwarz bounds their combined second moment by
$C_f(L_{ij}^2+L_{ji}^2)$.  Summing over unordered pairs proves the estimate.
\end{proof}

\begin{proposition}\label{prop:am-field-covariance}
For any finite list of smooth signals constant outside compact sets, and any
finite list of flag types, the laws of the independent fields
in~\eqref{eq:am-field} can be chosen so that, for every $g\in\mathcal K$,
the total variation distance between their joint law and its image under
$g$ is arbitrarily small on all sites at once.  With probability tending
to one as $N\to\infty$, every bit is one wherever its signal is at least
two, and zero wherever its signal is at most minus two.
\end{proposition}

\begin{proof}
We first consider one signal and one flag type.  The two noise matrices are
$T_0=I+B_N$ and $T_1=I+U_gB_NU_g^{-1}$, and the two means are $\mu_N$
and $U_g\mu_N$.  There is a finite set $J_N$ of $O(N^2)$ sites containing
the row and column supports of both perturbation matrices and the supports
of $\mu_N-q$ and $U_g\mu_N-q$, where $q$ is the constant value of
$\theta$ outside a compact set.  Outside $J_N$ both laws have exactly the same
independent background factor, namely the constant outside value of
$\theta$ plus independent noise with density $f^2$.  Thus their total
variation distance on the whole product space is exactly the total
variation distance of their restrictions to $J_N$.  This is the reason
that the calculation below proves an assertion about all sites, rather
than only about finitely many specified marginals.

On $J_N$ interpolate linearly between the matrices and means.  Write
$\Delta=T_1-T_0$, $T_t=T_0+t\Delta$, and $h=h_N$.
Each $T_t$ is the identity plus a positive semidefinite matrix, so
$\|T_t^{-1}\|_{\mathrm{op}}\le1$.  Also the resolvent identity gives
\begin{equation}\label{eq:am-resolvent}
 T_t^{-1}h=(I-tT_t^{-1}\Delta)T_0^{-1}h.
\end{equation}
Consequently
\[
 \sup_{0\le t\le1}\|T_t^{-1}h\|_2
 \le(1+\|\Delta\|_{\mathrm{op}})\|T_0^{-1}h\|_2,
 \qquad
 \|T_t^{-1}\Delta\|_{\HS}\le\|\Delta\|_{\HS}.
\]
Lemma~\ref{lem:am-density-speed}, integrated over $t$, bounds the
$L^2$ distance between the square-root densities by
\[
 C_f\bigl(\|\Delta\|_{\HS}
 +(1+\|\Delta\|_{\mathrm{op}})\|T_0^{-1}h\|_2\bigr).
\]
For probability densities $p,q$, Cauchy--Schwarz gives
$\frac12\|p-q\|_1\le\|\sqrt p-\sqrt q\|_2$.
Lemma~\ref{lem:am-estimates} therefore proves that the limiting total
variation error is at most $C(\eta+\sqrt\kappa)$.  Choose these two
parameters small and then $N$ large.  Thresholding the fields cannot
increase total variation.  For finitely many independent fields the joint
error is bounded by the sum of their errors, using independent couplings.

On the event $\|B_N\xi\|_\infty\le1/2$, the noise in each coordinate of
$Y_N$ has absolute value at most $3/2$.  Hence every signal value at least
two gives bit one, and every value at most minus two gives bit zero,
simultaneously.  The event has probability tending to one by
\eqref{eq:am-concentration}.  Outside the finite matrix support the same
statement holds deterministically because $|\xi_z|\le1$.
\end{proof}

\subsection{Finite barriers and exact covariance at chart edges}

For each family $j$ choose a nonzero tangent vector $v_j$ and a transverse
vector $w_j$, and use the flag $(v_j,w_j)$.  There will be two independent
bit fields of this type: the first selects subdivision points on the lines,
and the second selects which intervals between consecutive points become
barriers.  The constants and signals are chosen once from the fixed finite
chart data, before choosing the small smoothing parameters.

Call a line $\ell_j=c$, $c\in\OO$, a \emph{candidate line} if it meets the
interior of the square and
\begin{equation}\label{eq:am-candidate}
 |\bar c|\le4N.
\end{equation}
There are $O(N)$ candidates in each track and direction.  Indeed, their
ordinary levels lie in a fixed bounded interval, and their conjugate levels
lie in an interval of length $8N$, so~\eqref{eq:am-one-count} applies.
Let $Q$ bound $|\overline{\ell_j(u)}|$ for every translation vector in the
source and image charts, including the square corrections.

We next choose a conjugate-coordinate core that contains every endpoint
needed in comparing these candidates.  If a candidate $\ell_j=c$ meets a
nonparallel chart edge on $\ell_k=b$, then its intersection $p$ satisfies
\begin{equation}\label{eq:am-core-equations}
 \bar p=
 \begin{pmatrix}\bar\ell_j\\ \bar\ell_k\end{pmatrix}^{-1}
 \binom{\bar c}{\bar b}.
\end{equation}
Here $\bar\ell_j$ denotes the row of conjugate coefficients of $\ell_j$.
The inverse matrices exist: their determinants are the conjugates of
nonzero elements of $\mathbb Q(\sqrt5)$.  There are finitely many matrices,
and $b$ ranges over a fixed finite list.  Thus one may fix $R_0$ so that
$|\bar p|\le R_0N$ for every such intersection, both in source and image
charts and also for lines with $|\bar c|\le4N+Q$.  Include intersections
with square edges.  Increase $R_0$ by a fixed amount to include square
representatives of these points.  The intersection arithmetic ensures
$p\in L$.  A small conjugate determinant merely increases this fixed
radius; it causes no change in the $O(N^2)$ site count.

Choose $R_1>R_0+1$.  Take a smooth first signal $\theta_{1,j}$ equal to
$3$ on $\{|v|\le R_0\}$ and equal to $-3$ on
$\{|v|\ge R_1\}$.  The same radial choice would work for every $j$.
Take a smooth second signal $\theta_{2,j}$, constant $-3$ off a compact
set, such that
\begin{align}
 \theta_{2,j}(v)&=3
    &&\text{if } |v|\le R_1
          \text{ and }|\bar\ell_j(v)|\le1,
       \label{eq:am-positive-core}\\
 \theta_{2,j}(v)&=-3
    &&\text{if }|\bar\ell_j(v)|\ge3.
       \label{eq:am-negative-band}
\end{align}
For example choose smooth cutoffs $\alpha_j,\beta$ with
$\alpha_j=1$ on $\{|\bar\ell_j|\le1\}$ and zero on
$\{|\bar\ell_j|\ge3\}$, and with $\beta=1$ on the radius-$R_1$ ball
and compact support; then $6\alpha_j\beta-3$ has the stated properties.
This choice gives a negative band of fixed scaled width between any
possibly active line and the artificial cutoff $4N$.

Apply Proposition~\ref{prop:am-field-covariance} to these eight fields.
Call a sample \emph{successful} when all the positive and negative signal
requirements in that proposition hold simultaneously.  Its failure
probability $\alpha_N$ tends to zero for fixed small $\eta,\kappa$.
In a successful sample, every required chart-edge intersection is marked
by the first field, and every first mark lies in $|\bar z|<R_1N$.

Here is the partition algorithm.  On each candidate line orient its
intersection with the square by $v_j$.  Subdivide at all interior sites
whose first bit is one, and at the two geometric square endpoints.  Its
incoming square endpoint has the flag pointing into the segment, is a
site with that square representative, and is a forced first mark.  The
outgoing square endpoint is just a geometric endpoint.  Its positive
flag may be represented on the opposite side of the square, so it is not
used to start an interval in the current square.  For every interval of
the subdivision, use the second bit at its initial endpoint to decide
whether the closed interval is a barrier.  Add the square edges as barriers
regardless of the bits.  Lines lying along a square edge need
no further rule.  There are finitely many intervals, since every first
mark is among the $O(N^2)$ sites in a fixed conjugate ball.

The open connected components of the complement of the barriers give a
finite Boolean polygonal partition, denoted $\mathcal P_N$.  To justify
this assertion, take the finite arrangement of all the full lines supporting
the barrier segments, together with the square edges.  Each open chamber
is connected and avoids all segments, hence lies in one component.  Each
component is therefore, modulo empty interior, a union of finitely many
chambers.  Such a union belongs to the given Boolean algebra.  The
components are disjoint and cover the square modulo the discarded boundaries.
If the sample is unsuccessful, define $\mathcal P_N$ instead to be the
partition with one cell on each track.  Thus the algorithm is defined for
every sample, and its output always belongs to the countable set of finite
polygonal partitions.

\begin{lemma}\label{lem:am-partition-coupling}
The signals just chosen have the following property.  Given $\delta>0$,
one can choose $\eta,\kappa$ and then $N$ so that, for each
$g\in\mathcal K$, there is a coupling of two partitions $\mathcal P$ and
$\mathcal P'$ having the law of $\mathcal P_N$ for which, with probability
greater than $1-\delta$, every cell of $\mathcal P$ is contained in a
translation chart of $g$ and
\[
 \mathcal P'=g\mathcal P=\{gC:C\in\mathcal P\}.
\]
The equality is as Boolean partitions, including track labels.
\end{lemma}

\begin{proof}
For the eight bit fields, Proposition~\ref{prop:am-field-covariance}
provides a coupling of an original sample $b$ and a sample $b'$ with the
same law such that
\begin{equation}\label{eq:am-bit-coupling}
 b'(gz)=b(z)
\end{equation}
for both bits in every direction and every site, except on an event of
arbitrarily small probability $\varepsilon_N$.  Such a coupling follows
directly from the common part of the two probability measures: the mass
of their pointwise minimum is one minus their total variation distance;
on that mass make the samples identical, and couple the residual measures
arbitrarily.  This construction also applies to the countable product
space, since the measures here differ only on a finite coordinate set.
Intersect this coupling event with success of both samples.  The resulting
event has probability at least $1-\varepsilon_N-2\alpha_N$.
We prove the required equality deterministically on it.

First every fixed chart-edge supporting line meeting the square interior
is a full barrier in both samples for large $N$.  Its conjugate level is
bounded, so it is a candidate and satisfies the positive condition
\eqref{eq:am-positive-core}.  Every subdivision start lies in the first-mark
support and hence in the radius-$R_1$ ball after scaling.  Its second bit
is therefore one.  The incoming square endpoint is forced by the core,
and the final endpoint is supplied geometrically, so all of this line's
segment in the square is covered by activated intervals.  Square edges
are barriers by definition.  In particular the open cells of each sample
stay in the corresponding source or image continuity charts.

Fix one open source chart $U$, translated by $u$ onto the open image chart
$V$.  Consider a candidate line $L$ that meets $U$.  Its image is the
parallel allowable line $L+u$, with conjugate level shifted by
$\overline{\ell_j(u)}$.  First suppose that both lines are candidates.
Interior marked points of $L\cap U$ correspond exactly to interior marked
points of $(L+u)\cap V$ by~\eqref{eq:am-bit-coupling}.  At every transverse
crossing of $L$ with $\partial U$, the subdivision point was already
forced by~\eqref{eq:am-core-equations}; the corresponding intersection
with $\partial V$ is likewise forced independently.  Thus the subdivision
intervals in these open charts have exactly corresponding interiors and
endpoints under translation.

The direction of the flag matters at the endpoints.  If an interval
starts at $p$ and then enters $U$, its positive tangent flag selects $U$.
Its initial second bit consequently maps to the bit at $p+u$, which is
the start of the corresponding interval in $V$.  Hence the two intervals
have the same activation decision.  If the terminal endpoint is $q$, its
positive tangent flag may enter a different source chart and use a
translation different from $u$.  No identity between its second bit and
the bit at $q+u$ is needed: the preceding interval uses the bit at $p$.
The point $q+u$ is instead an independently forced destination subdivision
point when it lies on an interior chart edge, or a geometric truncation
when it is a square exit.  In particular, this comparison never inserts
a new random subdivision into an already activated interval.  Both source
and destination were subdivided at all transverse chart crossings before
the activation decisions were made.  Figure~\ref{fig:am-endpoints}
records the distinction.

\begin{figure}[tb]
\centering
\begin{tikzpicture}[x=1cm,y=1cm,>=Stealth,font=\small]
 \fill[blue!7] (0,2.0) rectangle (3.4,3.2);
 \fill[orange!10] (3.4,2.0) rectangle (6.5,3.2);
 \draw[gray] (3.4,2.0)--(3.4,3.2);
 \node at (1.7,3.0) {$U$};
 \node at (4.95,3.0) {$U'$};
 \draw[gray] (0.3,2.45)--(6.2,2.45);
 \draw[very thick,blue!70!black,->] (0.8,2.45)--(3.4,2.45);
 \fill[blue!70!black] (0.8,2.45) circle (2pt);
 \draw[fill=white,thick] (3.4,2.45) circle (2pt);
 \draw[orange!80!black,->,thick] (3.48,2.58)--(4.1,2.58);
 \node[below] at (0.8,2.4) {$p$};
 \node[below left] at (3.4,2.4) {$q$};
 \node[below,font=\scriptsize] at (5.0,2.38) {positive tangent flag};
 \fill[blue!7] (-0.4,-0.3) rectangle (3.0,0.9);
 \fill[orange!10] (3.9,-0.3) rectangle (7.0,0.9);
 \node at (1.3,0.65) {$U+u$};
 \node at (5.45,0.65) {$U'+u'$};
 \draw[very thick,blue!70!black,->] (0.4,0.15)--(3.0,0.15);
 \fill[blue!70!black] (0.4,0.15) circle (2pt);
 \draw[fill=white,thick] (3.0,0.15) circle (2pt);
 \fill[orange!80!black] (3.9,0.15) circle (2pt);
 \draw[orange!80!black,->,thick] (3.9,0.15)--(5.0,0.15);
 \node[below] at (0.4,0.1) {$p+u$};
 \node[below left] at (3.0,0.1) {$q+u$};
 \node[below right] at (3.9,0.1) {$q+u'$};
 \draw[->,dashed,gray] (3.2,1.95)--(3.0,0.95);
 \draw[->,dashed,orange!80!black] (3.7,1.95)--(3.9,0.95);
 \node[font=\scriptsize,anchor=east] at (2.8,1.4) {terminal endpoint};
 \node[font=\scriptsize,anchor=west] at (4.15,1.4) {flag image};
\end{tikzpicture}
\caption{Endpoint transport across a chart edge.  The interval starting at
$p$ uses the activation bit transported to $p+u$.  Its terminal endpoint
$q+u$ is forced by the destination core.  The positive flag based at $q$
uses the next chart translation $u'$; its image $q+u'$ need not coincide
with $q+u$.}
\label{fig:am-endpoints}
\end{figure}

At a vertex where several chart edges meet, an interval entering an open
chart still has its positive tangent flag in that chart.  If an interval
runs along a chart boundary, that boundary is already a full barrier and
requires no activation comparison.  Parallel lines either do not cross an
edge or coincide with its supporting line; these alternatives have just
been covered.  Intersections that only touch at a point make no difference
to Boolean cells.

It remains to handle the finite line cutoff.  If $L$ is a candidate but
$L+u$ is not, then
\[
 |\bar c|\ge4N-Q>3N
\]
for sufficiently large $N$.  Every second bit at a start on $L$ is zero
by~\eqref{eq:am-negative-band}, so the line contributes no active segment.
The reverse mismatch is handled by the same inequality on the destination
line.  A line absent because it misses the square cannot produce a
mismatch inside $U,V$, since a line meeting one of these corresponding
open charts translates to a line meeting the other.  Square seam
translations were included in the bound $Q$ and in the endpoint core.
Thus active barriers inside $U$ translate exactly to the active barriers
inside $V$ in every case.

Since all chart edges are full barriers in both samples, this equality
inside each pair of open charts identifies their connected complementary
components.  Each source cell therefore translates to precisely one
destination cell, and all destination cells occur.  It follows that
$\mathcal P'=g\mathcal P$.  Finally choose $\eta,\kappa$ small enough
and then $N$ large enough that
$\varepsilon_N+2\alpha_N<\delta$, simultaneously for the finite set
$\mathcal K$.
\end{proof}

The field construction and the endpoint argument have now produced a
partition law with the required covariance.  The remaining step is a
group-theoretic conversion.  It uses uniform matchings of equal shapes to
retain covariance, followed by an explicit passage from probability
measures to the finite sets in the definition of amenability.

\subsection{Shape matchings, subgroup measures, and finite F\o lner sets}

Two nonempty polygonal cells have the same \emph{shape} if one is a
translate of the other by an element of $\OO^2$ in ordinary square
coordinates; their tracks may differ.  The translation is unique.  Indeed,
a bounded set of positive area cannot be invariant modulo empty interior
under a nonzero translation: a linear functional increasing in that
direction has a finite essential supremum, which the translation would
increase.  A finite partition has a finite multiset of shapes, counting
multiplicities.

For every multiset of shapes that occurs for a partition of $mX$, fix one
baseline partition with that multiset.  There are only countably many
partitions, so these choices require no measurability qualification.
Given a sampled partition $\mathcal P$, choose uniformly a shape-preserving
bijection from its baseline partition onto $\mathcal P$.  On each baseline
cell use the unique translating identification to the assigned cell.
Their union is a piecewise translation $a\in F_m$.  Denote its probability
law by $\nu$.

In the successful coupling of Lemma~\ref{lem:am-partition-coupling}, the
partitions $\mathcal P$ and $g\mathcal P$ have the same multiset of
shapes, because each cell lies in one translation chart.  They consequently
have the same baseline.  Composition by $g$ bijects the shape-preserving
matchings to $\mathcal P$ with the shape-preserving matchings to
$g\mathcal P$.  Thus a uniform matching to the first, followed by $g$, is
a uniform matching to the second.  Extend this coupling by arbitrary
uniform matchings on the exceptional event.  Each marginal is still the
law prescribed above, and the resulting group elements are related by
left multiplication by $g$ on the successful event.  Hence
\begin{equation}\label{eq:am-reiter}
 \|g\nu-\nu\|_{\TV}<\delta
 \qquad(g\in\mathcal K).
\end{equation}
Here $(g\nu)(a)=\nu(g^{-1}a)$, and for measures on a countable set
$\|\sigma-\sigma'\|_{\TV}=\frac12\sum_a|\sigma(a)-\sigma'(a)|$.
For the concrete construction at a fixed scale $N$ there are in fact only
finitely many possible successful marked configurations and hence finitely
many partitions and matchings.  The following argument allows countable
support as well and proves subgroup inheritance at the same time.

\begin{lemma}\label{lem:am-reiter-folner}
Let $H$ be a subgroup of a countable group $J$.  Suppose that for every
finite $K\subset H$ and every $\varepsilon>0$ there is a probability
measure $\nu$ on $J$ such that
\[
 \sum_{k\in K}\|k\nu-\nu\|_1<\varepsilon.
\]
Then for every finite $K\subset H$ and every $\varepsilon>0$ there is a
nonempty finite $A\subset H$ such that
\[
 \sum_{k\in K}|kA\mathbin\triangle A|<\varepsilon|A|.
\]
\end{lemma}

\begin{proof}
Choose one representative $t$ in every right coset $Ht$ of $H$ in $J$,
and write each $x\in J$ uniquely as $x=h t$.  Define $\pi(x)=h$.
Then $\pi(kx)=k\pi(x)$ for $k\in H$.  The pushforward
$p=\pi_*\nu$ is a probability measure on $H$, and summing over fibers
shows
\[
 \|kp-p\|_1\le\|k\nu-\nu\|_1.
\]
Start with the sum on the right less than $\varepsilon/2$.  If $K$ is
empty any singleton suffices, so assume $K$ is nonempty.  Choose a finite
$E\subset H$ with $p(E)>1-\varepsilon/(8|K|)$, decreasing this error
further if necessary to ensure $p(E)>0$, and put
$q=p\mathbf1_E/p(E)$.  Direct calculation gives
$\|q-p\|_1=2(1-p(E))$.  Therefore
\[
 \sum_{k\in K}\|kq-q\|_1
 \le \sum_{k\in K}\|kp-p\|_1
       +4|K|(1-p(E))<\varepsilon.
\]
For $t>0$ let $A_t=\{h\in H:q(h)>t\}$.  These sets are finite, and the
identities for real numbers $a,b\ge0$,
\[
 \int_0^\infty\mathbf1_{\{a>t\}}\,dt=a,
 \qquad
 \int_0^\infty
 |\mathbf1_{\{a>t\}}-\mathbf1_{\{b>t\}}|\,dt=|a-b|,
\]
give, after finite summation,
\[
 \int_0^\infty |A_t|\,dt=1,
 \qquad
 \int_0^\infty\sum_{k\in K}|kA_t\mathbin\triangle A_t|\,dt
 =\sum_{k\in K}\|kq-q\|_1<\varepsilon.
\]
If every nonempty $A_t$ had its summed boundary at least
$\varepsilon|A_t|$, integration would contradict these two formulas.
For some $t$ the nonempty finite set $A=A_t$ therefore satisfies the
asserted strict inequality.
\end{proof}

\begin{proof}[Proof of Theorem~\ref{thm:amenability}]
Fix a subgroup $H\le F_m$, a finite $K\subset H$, and $\varepsilon>0$.
For nonempty $K$, apply the construction with $\mathcal K=K\cup K^{-1}$
and with the total variation error in~\eqref{eq:am-reiter} less than
$\varepsilon/(4|K|)$.  It supplies a probability measure on $F_m$ whose
summed $\ell^1$ error for $K$ is less than $\varepsilon/2$.
Lemma~\ref{lem:am-reiter-folner} supplies a nonempty finite $A\subset H$
with summed boundary less than $\varepsilon|A|$, and hence
$|kA\mathbin\triangle A|<\varepsilon|A|$ for every $k\in K$.
The empty $K$ case is immediate.  Taking $H=F_m$ proves amenability of
$F_m$, and taking $H=A_m$ gives the assertion needed for the simple group.
All choices were made with $m$ fixed; no passage to infinitely many tracks
has occurred.
\end{proof}
